\section{Resultados de la línea base}
\label{sec:resultados}

\guia{Completar la tabla con mediciones reales de cada modelo y nivel.
$N$ es el número de problemas, $n$ las muestras por problema y $k$ el
número de intentos de la métrica. Si hay varios valores de $k$, añadir
filas o columnas identificadas. Los campos vacíos no representan cero.}

\begin{table}[htbp]
  \centering\small
  \caption{Resultados por modelo y nivel de dificultad.}
  \label{tab:lineabase}
  \begin{tabularx}{\linewidth}{@{}Ylccccc@{}}
    \toprule
    Modelo & Dificultad & $N$ & $n$ & $k$ & \passuno & \passk \\
    \midrule
    Objetivo & Básica & \valor & \valor & \valor & \valor & \valor \\
    Objetivo & Intermedia & \valor & \valor & \valor & \valor & \valor \\
    Objetivo & Alta & \valor & \valor & \valor & \valor & \valor \\
    \midrule
    Comparación 1 & Básica & \valor & \valor & \valor & \valor & \valor \\
    Comparación 1 & Intermedia & \valor & \valor & \valor & \valor & \valor \\
    Comparación 1 & Alta & \valor & \valor & \valor & \valor & \valor \\
    \midrule
    Comparación 2 & Básica & \valor & \valor & \valor & \valor & \valor \\
    Comparación 2 & Intermedia & \valor & \valor & \valor & \valor & \valor \\
    Comparación 2 & Alta & \valor & \valor & \valor & \valor & \valor \\
    \bottomrule
  \end{tabularx}
\end{table}

\pendiente{comparar modelos y niveles, describir el efecto de aumentar
los intentos e identificar patrones que requieran explicación.}

\guia{Indicar la incertidumbre de las mediciones cuando se disponga
de ella. Si se informa una métrica global, conservar el desglose por
dificultad solicitado por el curso.}
